\begin{figure}[t]\centering\includegraphics[width=0.9\columnwidth]{tradeoff.pdf}\caption{Skill against climate at $c{=}10$ for every configuration run at the training forcing $F{=}20$: main systems and the noise, MLP-input and degree ablations (mean over seeds 0--2 for each point; log axis for spec\_err). Colours: grey no closure, blue polynomial, orange MLP, red AR(1) variants.}\label{fig:tradeoff}\end{figure}
\begin{table}[t]\centering\caption{H5 and noise amplitude, $c{=}10$, seeds 0--2 (mean $\pm$ std). The deterministic polynomial and the fitted AR(1) closure are the main-group rows of Table~\ref{tab:main} (five seeds).}\label{tab:ar1}
\footnotesize\resizebox{\columnwidth}{!}{\begin{tabular}{lcccc}
\toprule
Closure noise & valid\_time & var\_ratio & w1\_pdf & spec\_err \\
\midrule
AR(1) $\sigma\times0.5$ & 1.367 $\pm$ 0.016 & 0.992 $\pm$ 0.002 & 0.078 $\pm$ 0.006 & 0.053 $\pm$ 0.002 \\
AR(1) $\sigma\times1.5$ & 1.129 $\pm$ 0.014 & 1.004 $\pm$ 0.002 & 0.104 $\pm$ 0.005 & 0.071 $\pm$ 0.003 \\
White noise ($\phi\!=\!0$) & 1.399 $\pm$ 0.025 & 0.985 $\pm$ 0.002 & 0.101 $\pm$ 0.009 & 0.067 $\pm$ 0.005 \\
\bottomrule
\end{tabular}
}
\end{table}
\begin{table}[t]\centering\caption{Polynomial degree, $c{=}10$, seeds 0--2. Degree 4 is the main-group polynomial of Table~\ref{tab:main} (five seeds).}\label{tab:deg}
\footnotesize\resizebox{\columnwidth}{!}{\begin{tabular}{lccccc}
\toprule
Polynomial & offline\_r2 & valid\_time & mean\_err & w1\_pdf & spec\_err \\
\midrule
Degree 1 & 0.791 $\pm$ 0.001 & 1.199 $\pm$ 0.045 & 0.192 $\pm$ 0.007 & 0.197 $\pm$ 0.006 & 0.099 $\pm$ 0.005 \\
Degree 2 & 0.830 $\pm$ 0.001 & 1.319 $\pm$ 0.017 & 0.598 $\pm$ 0.029 & 0.607 $\pm$ 0.029 & 0.381 $\pm$ 0.026 \\
Degree 3 & 0.841 $\pm$ 0.001 & 1.340 $\pm$ 0.032 & 0.315 $\pm$ 0.094 & 0.321 $\pm$ 0.094 & 0.191 $\pm$ 0.057 \\
Degree 4 & 0.847 $\pm$ 0.001 & 1.372 $\pm$ 0.029 & 0.222 $\pm$ 0.043 & 0.229 $\pm$ 0.042 & 0.137 $\pm$ 0.023 \\
Degree 5 & 0.847 $\pm$ 0.001 & 1.380 $\pm$ 0.032 & 0.167 $\pm$ 0.016 & 0.174 $\pm$ 0.016 & 0.105 $\pm$ 0.006 \\
Degree 6 & 0.848 $\pm$ 0.001 & 1.394 $\pm$ 0.036 & 0.179 $\pm$ 0.033 & 0.188 $\pm$ 0.033 & 0.109 $\pm$ 0.018 \\
\bottomrule
\end{tabular}
}
\end{table}
\begin{table}[t]\centering\caption{MLP input width, $c{=}10$, seeds 0--2 (hidden width 32 in all rows). The local-input MLP is the main-group row of Table~\ref{tab:main} (five seeds).}\label{tab:mlpin}
\footnotesize\resizebox{\columnwidth}{!}{\begin{tabular}{lccccc}
\toprule
MLP input & offline\_r2 & valid\_time & mean\_err & w1\_pdf & spec\_err \\
\midrule
$X_{k-1},X_k,X_{k+1}$ & 0.878 $\pm$ 0.002 & 1.398 $\pm$ 0.013 & 0.019 $\pm$ 0.003 & 0.056 $\pm$ 0.004 & 0.031 $\pm$ 0.004 \\
$X_{k-3..k+3}$ & 0.923 $\pm$ 0.001 & 1.538 $\pm$ 0.012 & 0.053 $\pm$ 0.049 & 0.062 $\pm$ 0.043 & 0.043 $\pm$ 0.024 \\
\bottomrule
\end{tabular}
}
\end{table}
\begin{table*}[t]\centering\caption{Forcing shift: closures trained at $F{=}20$, slow model and truth at the deployment $F$ ($c{=}10$, seeds 0--2). The values at the training forcing $F{=}20$ are in Table~\ref{tab:main} (five seeds).}\label{tab:shift}
\footnotesize\resizebox{0.85\textwidth}{!}{\begin{tabular}{llccccc}
\toprule
Deploy $F$ & System & valid\_time & mean\_err & var\_err & w1\_pdf & spec\_err \\
\midrule
20 (trained) & Polynomial (deg 4) & 1.372 $\pm$ 0.029 & 0.222 $\pm$ 0.043 & 0.039 $\pm$ 0.008 & 0.229 $\pm$ 0.042 & 0.137 $\pm$ 0.023 \\
20 (trained) & MLP (32x2) & 1.403 $\pm$ 0.037 & 0.209 $\pm$ 0.028 & 0.036 $\pm$ 0.005 & 0.217 $\pm$ 0.027 & 0.124 $\pm$ 0.014 \\
20 (trained) & AR(1) stochastic & 1.250 $\pm$ 0.027 & 0.061 $\pm$ 0.006 & 0.002 $\pm$ 0.001 & 0.083 $\pm$ 0.004 & 0.057 $\pm$ 0.002 \\
\midrule
18 & Polynomial (deg 4) & 1.485 $\pm$ 0.043 & 0.014 $\pm$ 0.004 & 0.003 $\pm$ 0.001 & 0.070 $\pm$ 0.003 & 0.039 $\pm$ 0.003 \\
18 & MLP (32x2) & 1.538 $\pm$ 0.050 & 0.019 $\pm$ 0.002 & 0.002 $\pm$ 0.000 & 0.057 $\pm$ 0.001 & 0.038 $\pm$ 0.000 \\
18 & AR(1) stochastic & 1.367 $\pm$ 0.031 & 0.089 $\pm$ 0.003 & 0.004 $\pm$ 0.001 & 0.107 $\pm$ 0.006 & 0.078 $\pm$ 0.006 \\
\midrule
22 & Polynomial (deg 4) & 1.357 $\pm$ 0.044 & 0.036 $\pm$ 0.002 & 0.010 $\pm$ 0.001 & 0.039 $\pm$ 0.002 & 0.010 $\pm$ 0.001 \\
22 & MLP (32x2) & 1.394 $\pm$ 0.037 & 0.003 $\pm$ 0.004 & 0.004 $\pm$ 0.001 & 0.034 $\pm$ 0.004 & 0.006 $\pm$ 0.001 \\
22 & AR(1) stochastic & 1.211 $\pm$ 0.062 & 0.033 $\pm$ 0.002 & 0.017 $\pm$ 0.001 & 0.046 $\pm$ 0.002 & 0.015 $\pm$ 0.000 \\
\bottomrule
\end{tabular}
}
\end{table*}

\textbf{H5 and noise amplitude (Table~\ref{tab:ar1}).} The ablation rows use seeds 0--2 and are set against the five-seed main-group values of Table~\ref{tab:main}, so these comparisons are not seed-matched. White noise of the same variance ($\phi{=}0$) gives spec\_err 0.067 against \rhval{main/ar-1-stochastic/f20_c10/spec_err/mean} for the fitted AR(1) ($\phi$ fitted to 0.982 for seed 0), and both are far below the deterministic closure (\rhval{main/polynomial-deg-4/f20_c10/spec_err/mean}); white noise retains a longer valid time (1.399 against \rhval{main/ar-1-stochastic/f20_c10/valid_time/mean}). The ordering in H5 holds in the means, but the gap is small compared with the gap to the deterministic closure, so in these means most of the reduction in spec\_err is already obtained with white noise of the fitted variance, and we cannot attribute the remainder to temporal correlation without more seeds or a test; the white-noise run also has a lower variance ratio (0.985 against \rhval{main/ar-1-stochastic/f20_c10/var_ratio/mean}). Amplitude trades skill for climate: valid time falls monotonically from \rhval{main/polynomial-deg-4/f20_c10/valid_time/mean} (no noise) through 1.367 ($\times0.5$) and \rhval{main/ar-1-stochastic/f20_c10/valid_time/mean} to 1.129 ($\times1.5$), while spec\_err is lowest at $\times0.5$ (0.053) and similar at $\times1$ (\rhval{main/ar-1-stochastic/f20_c10/spec_err/mean}) and worse at $\times1.5$ (0.071). The fitted amplitude is therefore not obviously optimal: half of it gave a similar climate and a valid time close to that of the deterministic closure (3 seeds against 5, no test).

\textbf{Polynomial degree (Table~\ref{tab:deg}), a failure case.} Over the ablation degrees 1, 2, 3, 5 and 6, offline $R^2$ never decreases with degree (0.791 to 0.848), and valid time increases with it (1.199 to 1.394), but the climate does not: degree 2 has by far the worst climate (spec\_err 0.381, mean\_err 0.598), worse than the linear closure (0.099), while degree 3 is at 0.191, degrees 5 and 6 are at 0.105 and 0.109, and degree 4 (main group, five seeds) is at \rhval{main/polynomial-deg-4/f20_c10/spec_err/mean}. We have no isolating experiment for why degree 2 fails and do not claim a mechanism; the point is that offline $R^2$ and forecast skill did not predict the climate of the coupled model.

\textbf{MLP input width (Table~\ref{tab:mlpin}).} Giving the MLP the neighbouring slow variables lifts offline $R^2$ from \rhval{main/mlp-32x2/f20_c10/offline_r2/mean} (local input, main group, five seeds) to 0.878 (three inputs) and 0.923 (seven inputs). Climate improves at width three (spec\_err \rhval{main/mlp-32x2/f20_c10/spec_err/mean} to 0.031, mean\_err \rhval{main/mlp-32x2/f20_c10/mean_err/mean} to 0.019) while valid time does not (1.398 against \rhval{main/mlp-32x2/f20_c10/valid_time/mean}); the seven-input net had the best valid time (1.538) but a large seed-to-seed spread of mean\_err (0.053 $\pm$ 0.049). This exploratory result (3 seeds, not seed-matched to the main group) suggests that the local-input restriction, rather than the function class, may be what limited the MLP in the main comparison; it was not a registered hypothesis and the width was not chosen on a separate split.

\textbf{Forcing shift (Table~\ref{tab:shift}).} No closure diverged. At $F{=}18$ the AR(1) closure is worse than the deterministic ones on spec\_err (0.078 against 0.039 for the polynomial) and w1\_pdf (0.107 against 0.070), whereas at $F{=}20$ it was better; its noise amplitude was fitted at $F{=}20$ and kept fixed, which may explain it (not isolated by a run). At $F{=}22$ AR(1) is again behind the polynomial on spec\_err (0.015 against 0.010) and w1\_pdf (0.046 against 0.039) and has the largest var\_err (0.017), while the MLP is best on all climate columns of the table.
